\documentclass[11pt,a4paper]{article}
\usepackage[margin=25mm]{geometry}
\usepackage{amsmath,amssymb,amsthm,booktabs}
\newtheorem{lemma}{Lemma}
\newcommand{\Aut}{\operatorname{Aut}}
\newcommand{\C}{\operatorname{C}}
\title{Inner reflexibility when one rotation has order three}
\author{Human Compression project}\date{}
\begin{document}\maketitle
Let $G=O'N$ and $A=\Aut(G)=G{:}2$. We use the fixed source's convention:
ordered generating pairs $(R,S)$ with $|R|=3$, $|S|=q>2$ and $|RS|=2$
are counted up to $A$. There is no division by duality. The accepted
inversion correspondence identifies those with an inner simultaneous
inverter with triples $(a,b,c)$ of inner involutions satisfying
\[
 |ab|=3,\qquad |bc|=q,\qquad [a,c]=1,\qquad a\ne c,
 \qquad \langle ab,bc\rangle=G.
\]
The correspondence is $R=ab$, $S=bc$; its inverse is unique for a
generating pair. In particular, $A$ acts freely on these triples.

\section{The two possible partners}
Fix an inner involution $a$, and put $J=\C_G(a)$ and $K=\C_A(a)$.
The previously accepted character calculation gives one inner involution
class and the exact values
\[
 |J|=161280,\qquad |K|=322560,\qquad
 \kappa_3=20608,\qquad \kappa_2=1750.
\]
Here $\kappa_3$ counts inner involutions $b$ with $|ab|=3$, and
$\kappa_2$ counts inner involutions $c\ne a$ commuting with $a$.
These are inputs from the earlier character calculation, rather than
counts inferred from the new searches.

The retained standard-generator program gives an inner element $u$ of
order $16$ and an outer involution $v$. Literal multiplication checks
that $a=u^8$ commutes with both. Consequently
\[
 J_0=\langle u,u^v\rangle\le J.
\]
The old variable names in that program are not used as group
identifications. The full words, their inner or outer parity, and the
displayed relations are checked in the accepted faithful permutation
model of $A$ of degree $245520$.

The color argument used for the outer row has the following useful
form. Label the $J_0$-orbits on the full permutation domain $X$ by
distinct colors. On one $J_0$-orbit $\Omega$, color $x$ by the orbit
containing $x^b$.

\begin{lemma}\label{saturation}
Suppose that $|\Omega|=6720$ and this coloring has an orbit of size
$20160$ under the image of $J_0$. If its stabilizer contains a displayed
subgroup of order $8$, then $J_0=J$, the action on $\Omega$ is faithful,
and the coloring stabilizer has order $8$. If the displayed generators
lift to elements commuting with $b$, then $|\C_J(b)|=8$.
\end{lemma}
\begin{proof}
The image of $J_0$ has order at least $20160\cdot8=161280$.
Its order is at most $|J_0|$, which is at most $|J|=161280$.
Equality holds in both inequalities, proving the group equality and
faithfulness. Orbit--stabilizer also shows that the displayed subgroup
is the whole coloring stabilizer. An element of $J$ centralizing $b$
preserves the coloring: it permutes $\Omega$ and preserves every
$J$-orbit on $X$. Conversely, the specified lifts of the full coloring
stabilizer centralize $b$. Faithfulness gives the last conclusion.
\end{proof}

For a second retained partner, the same faithful action gives a color
orbit of size $448$ and a displayed stabilizer of order $360$. Their
product is again $161280$, and the lifts again centralize $b$. Thus
the second common centralizer has order $360$.

The two partners have different color histograms. This proves that they
are not conjugate under $J$: for $g\in J$, the color assigned by $b^g$
at $x$ is the color assigned by $b$ at $x^{g^{-1}}$, because the final
action of $g$ preserves every $J$-orbit. Conjugation therefore permutes
the positions without changing their color multiplicities. The two
partner orbits have sizes $20160$ and $448$, whose sum is exactly
$\kappa_3$. Hence they exhaust all possible inner partners.

One partner is an explicit conjugate of $a$ by a retained word. The
second is obtained from a pair $a',b'\in J_0$ and a literal transport
$a'^h=a$. The transport uses the elementary identity
\[
 x^{(xy)^{(k+1)/2}}=y
 \quad\text{when}\quad x^2=y^2=1,\quad |xy|=k\text{ is odd}.
\]
It follows by writing $z=xy$ and using $xzx=z^{-1}$.
Every product order and the final transport equality is checked as a
permutation identity. The first exploratory transport used the opposite
exponent and failed its equality check; that failed execution remains
preserved. It is not a premise of the corrected certificate.

\section{A complete finite domain for the third involution}
Write $H_1=\C_J(b_1)$ and $H_2=\C_J(b_2)$, of orders $8$ and $360$.
In the faithful action of degree $6720$, the reader forms disjoint
conjugacy orbits of involutions with total size $1751$, including $a$.
For each $H_i$, it verifies a disjoint orbit partition of the remaining
$1750$ elements. Every orbit representative has an explicit word in
$u,u^v$; its full lift is checked to be an involution other than $a$,
to commute with $a$, and to restrict to the displayed permutation.
The entire orbit then lies in $J$, since $H_i\le J$.

This last membership check is necessary. An arbitrary involution of the
surrounding symmetric group need not belong to $J$. The word lifts and
the complete $H_i$-cover prove membership of every element actually used.
The independent upper bound $\kappa_2=1750$ now proves that the covered
set is the full required domain. No exhaustion claim returned by a
conjugacy-class algorithm is needed as a final premise.

The returned partitions have $250$ and $18$ orbits. For each
representative $c$, the reader verifies $|b_i c|$ directly. The
generation reader then forms $R=ab_i$ and $S=b_i c$ and checks their
point orbit in the full model. The actual socle orbit has size $122760$.
The complete acquired proper-maximal orders are
\[
 3753792^2,\quad175560,\quad161280,\quad25920^2,\quad14880^2,
 \quad10752,\quad7920^2,\quad2520^2,
\]
where superscripts indicate multiplicities. None is divisible by
$122760$. Thus $\langle R,S\rangle=G$ exactly when its orbit of a point
in this socle orbit has size $122760$: a proper subgroup is contained
in a proper maximal subgroup, and orbit--stabilizer would force the
forbidden divisibility. Completeness of the maximal list is the
explicit borrowed classification premise already used for the outer row.

\section{The full automorphism denominator}
For $i=1,2$, let $M_i(q)$ be the sum of the sizes of those $H_i$-orbits
whose representatives give generating pairs of type $(3,q)$.
The group $K$ normalizes $J$, so it permutes the two $J$-orbits of
partners, preserving their sizes. Since their sizes differ, each orbit
is separately $K$-invariant. Consequently
\[
 |\C_A(a,b_i)|=\frac{|K|}{|J|/|H_i|}=2|H_i|.
\]
This group acts freely on the generating choices of $c$.
Generation and product order are constant on each $H_i$-orbit, hence
the number of $A$-classes with an inner inverter is
\[
 \boxed{\displaystyle I(3,q)=\frac{M_1(q)}{16}
                         +\frac{M_2(q)}{720}.}
\]
No explicit fusion of $H_i$-orbits under the larger stabilizer is
required: its free action on their union determines the quotient
cardinality. In particular, each $M_i(q)$ is divisible by $2|H_i|$.
Dividing only by $|H_i|$ would incorrectly count the source's
full-automorphism classes twice.

The generation replay supplies the masses, and the chiral count is
\[
 C(3,q)=T(3,q)-I(3,q)-O(3,q),
\]
using the already proved total $T$ and outer count $O$. The uniqueness
of the inverter makes the inner and outer sets disjoint. The source's
asserted aggregate inner and chiral counts are not inputs to this
calculation. Both literal replays have returned; the exact results are
shown below.

\begin{center}\begin{tabular}{rrrrr}\toprule
$q$&Total&Inner&Outer&Chiral\\\midrule
3&0&0&0&0\\
4&0&0&0&0\\
5&0&0&0&0\\
6&0&0&0&0\\
7&0&0&0&0\\
8&10&2&0&8\\
10&7&0&1&6\\
11&37&3&2&32\\
12&0&0&0&0\\
14&10&1&3&6\\
15&6&5&1&0\\
16&68&10&0&58\\
19&57&14&3&40\\
20&44&0&2&42\\
28&20&4&2&14\\
31&25&0&1&24\\
\midrule
Total&284&39&15&230\\
\bottomrule\end{tabular}\end{center}


The first partner class contributes all $39$ inner classes. None of the
representatives in the second partner class generates $G$. Together with the independently proved outer total $15$, the
accepted generating-pair total $284$ gives $230$ chiral classes. These
recover the aggregate counts in Section 7 on printed page 219.

For $(3,15)$, all inner generating representatives occur in the first
partner class and their total weight is $80$. Thus there are $80/16=5$
inner classes. The accepted outer count is one, so all six generating
classes are reflexible. This proves the claimed absence of chirality
for this type but corrects the split in the proof of Lemma 5.1 on printed
page 218: five inner and one outer replace four inner and two outer.
The same table proves eight chiral classes of type $(3,8)$ and six of
type $(3,10)$, and confirms the absence of an inner inverter in the
latter type. These are separate source assertions, not new conjectures.

The inversion correspondence, character coefficients, faithful named
model and maximal-subgroup classification are attributed imported
inputs. Orbit--stabilizer, the color construction and the transport
identity supply the elementary argument. The contribution of this
application is the explicit finite certificate for the complete inner
order-three row; no priority claim for these general ideas is made.
Whole-source completion, other rotation rows, and Lean verification
remain separate from this component.
\end{document}
